\documentclass[11pt,a4paper]{article}

\newcommand{\doctitle}{PRS XVA Model (MKM-XV-001)}
\newcommand{\docsubtitle}{Credit, Funding and Capital Valuation Adjustments on the Physical Risk Swap}
\newcommand{\docversion}{0.9}
\newcommand{\docdate}{9-October-2026}
\newcommand{\docauthor}{David K Kelly}

% MKM Research Labs — Shared Document Header
% Usage: % MKM Research Labs — Shared Document Header
% Usage: % MKM Research Labs — Shared Document Header
% Usage: \input{../shared/mkm_header}  (before \begin{document})
%
% Required commands before \input:
%   \newcommand{\doctitle}{...}
%   \newcommand{\docsubtitle}{...}
%   \newcommand{\docversion}{1.0}
%   \newcommand{\docdate}{13-February-2026}
%   \newcommand{\docauthor}{David K Kelly}

% ─── Packages ───────────────────────────────────────────────────────────────
\usepackage[utf8]{inputenc}
\usepackage{amsmath,amsfonts,amssymb}
\usepackage{listings}
\usepackage{color,xcolor}
\usepackage{graphicx}
\usepackage{hyperref}
\usepackage[toc,page]{appendix}
\usepackage{tabto}
\usepackage{float}
\usepackage{comment}
\usepackage[en-GB]{datetime2}
\usepackage{setspace}
\usepackage{subcaption}
\usepackage{booktabs}
\usepackage{longtable}
\usepackage{adjustbox}
\usepackage{enumitem}
\usepackage{fancyhdr}
\usepackage{lastpage}

% ─── Page geometry ──────────────────────────────────────────────────────────
\onehalfspacing
\usepackage{geometry}
\geometry{a4paper, total={170mm,247mm}, left=20mm, top=25mm, bottom=25mm}

% ─── Headers and footers ───────────────────────────────────────────────────
\pagestyle{fancy}
\fancyhf{}
\fancyhead[L]{\small\textcolor{gray}{MKM Research Labs}}
\fancyhead[R]{\small\textcolor{gray}{\doctitle{} --- v\docversion}}
\fancyfoot[C]{\small\textcolor{gray}{CONFIDENTIAL}}
\fancyfoot[R]{\small\textcolor{gray}{Page \thepage\ of \pageref{LastPage}}}
\renewcommand{\headrulewidth}{0.4pt}
\renewcommand{\footrulewidth}{0.4pt}

% ─── Colours ────────────────────────────────────────────────────────────────
\definecolor{mkmblue}{HTML}{1565C0}
\definecolor{codegreen}{rgb}{0,0.6,0}
\definecolor{codegray}{rgb}{0.5,0.5,0.5}
\definecolor{codepurple}{rgb}{0.58,0,0.82}
\definecolor{backcolour}{rgb}{0.95,0.95,0.92}

% ─── Hyperref setup ─────────────────────────────────────────────────────────
\hypersetup{
    colorlinks=true,
    linkcolor=mkmblue,
    citecolor=mkmblue,
    urlcolor=mkmblue
}

% ─── Code listings ──────────────────────────────────────────────────────────
\lstdefinestyle{mystyle}{
    backgroundcolor=\color{backcolour},
    commentstyle=\color{codegreen},
    keywordstyle=\color{magenta},
    numberstyle=\tiny\color{codegray},
    stringstyle=\color{codepurple},
    basicstyle=\ttfamily\footnotesize,
    breakatwhitespace=false,
    breaklines=true,
    captionpos=b,
    keepspaces=true,
    numbers=left,
    numbersep=5pt,
    showspaces=false,
    showstringspaces=false,
    showtabs=false,
    tabsize=2
}
\lstset{style=mystyle}

% ─── Section formatting ────────────────────────────────────────────────────
\setcounter{tocdepth}{3}
\setcounter{secnumdepth}{3}

% ─── Shared macros ──────────────────────────────────────────────────────────
\newcommand{\mkmtitlepage}{%
    \begin{titlepage}
    \begin{center}
    \vspace*{1cm}
    \begin{figure}[H]\centering
        \includegraphics[width=0.4\textwidth]{../shared/Logos/MKM_logo.png}
    \end{figure}
    \vspace{1.5cm}
    {\Huge\bfseries\textcolor{mkmblue}{\doctitle}\par}
    \vspace{0.5cm}
    {\Large\textcolor{gray}{\docsubtitle}\par}
    \vspace{2cm}
    {\large MKM Research Labs\par}
    \vspace{0.5cm}
    {\large Version \docversion{} --- \docdate\par}
    \vspace{0.5cm}
    {\large \docauthor\par}
    \vspace{2cm}
    \textcolor{red}{\textbf{CONFIDENTIAL --- PROPRIETARY}}
    \end{center}
    \end{titlepage}
}

\newcommand{\mkmlegalpage}{%
    \clearpage
    \thispagestyle{empty}
    \begin{center}
    \large\textbf{Legal Notice}
    \vspace{1em}

    \textbf{PROPRIETARY AND CONFIDENTIAL}
    \end{center}

    \noindent This document, including all algorithms, models, methodology, formulas,
    calculations, and intellectual concepts contained herein, constitutes the exclusive
    intellectual property and confidential trade secrets of MKM Research Labs (``MKM'').

    \vspace{0.5em}
    \noindent Any usage, reproduction, distribution, modification, or disclosure of this
    document or any portion thereof, without the express prior written authorisation from
    MKM Research Labs is strictly prohibited and constitutes an infringement of intellectual
    property rights.

    \vspace{0.5em}
    \noindent The document is provided on an ``as is'' basis. MKM Research Labs makes no
    representations or warranties of any kind, express or implied, regarding its accuracy,
    reliability, or suitability for any particular purpose.

    \vspace{1em}
    \begin{center}
    All rights reserved. \copyright\ 2021--2026 MKM Research Labs.

    \vspace{0.5em}
    Governed by the laws of the United Kingdom.
    \end{center}
    \clearpage
}

\newcommand{\mkmhistorypage}[1]{%
    \section{Document History}
    \begin{table}[ht]
        \centering
        \begin{tabular}{l|l|l|l}
            \textbf{Date} & \textbf{Version} & \textbf{Description} & \textbf{Author} \\
            \hline
            #1
        \end{tabular}
        \label{tab:revision_history}
    \end{table}
    \clearpage
}

% ─── Bank validation protocol section commands ──────────────────────────────
% Use these to ensure every model doc has the required sections.
% They are just styled \section calls — content is provided by each doc.

\newcommand{\mkmexecsummary}{\section{Executive Summary}}
\newcommand{\mkmpurpose}{\section{Model Purpose and Scope}}
\newcommand{\mkmframework}{\section{Mathematical Framework}}
\newcommand{\mkminputs}{\section{Input Parameters and Calibration}}
\newcommand{\mkmimplementation}{\section{Implementation Details}}
\newcommand{\mkmsensitivity}{\section{Sensitivity Analysis}}
\newcommand{\mkmvalidation}{\section{Validation and Backtesting}}
\newcommand{\mkmtestresults}{\subsection{Automated Test Results}\input{test_results}}
\newcommand{\mkmlimitations}{\section{Model Limitations and Known Weaknesses}}

\newcommand{\mkmgovernance}{\section{Governance Validation Questions}
\noindent The following nine questions are mandated by Chapter~5 of the
\emph{Model Risk Governance for Vendors} handbook and must be explicitly
addressed for every model in the inventory. Response status is tracked in the
Model Governance Dashboard (MRC portal). \par\medskip}

\newcommand{\mkmvqitem}[3]{\subsection{Q#1 --- #2}
\noindent\textit{``#3''}\par\medskip\noindent}

\newcommand{\mkmchangelog}{\section{Change History}}
  (before \begin{document})
%
% Required commands before \input:
%   \newcommand{\doctitle}{...}
%   \newcommand{\docsubtitle}{...}
%   \newcommand{\docversion}{1.0}
%   \newcommand{\docdate}{13-February-2026}
%   \newcommand{\docauthor}{David K Kelly}

% ─── Packages ───────────────────────────────────────────────────────────────
\usepackage[utf8]{inputenc}
\usepackage{amsmath,amsfonts,amssymb}
\usepackage{listings}
\usepackage{color,xcolor}
\usepackage{graphicx}
\usepackage{hyperref}
\usepackage[toc,page]{appendix}
\usepackage{tabto}
\usepackage{float}
\usepackage{comment}
\usepackage[en-GB]{datetime2}
\usepackage{setspace}
\usepackage{subcaption}
\usepackage{booktabs}
\usepackage{longtable}
\usepackage{adjustbox}
\usepackage{enumitem}
\usepackage{fancyhdr}
\usepackage{lastpage}

% ─── Page geometry ──────────────────────────────────────────────────────────
\onehalfspacing
\usepackage{geometry}
\geometry{a4paper, total={170mm,247mm}, left=20mm, top=25mm, bottom=25mm}

% ─── Headers and footers ───────────────────────────────────────────────────
\pagestyle{fancy}
\fancyhf{}
\fancyhead[L]{\small\textcolor{gray}{MKM Research Labs}}
\fancyhead[R]{\small\textcolor{gray}{\doctitle{} --- v\docversion}}
\fancyfoot[C]{\small\textcolor{gray}{CONFIDENTIAL}}
\fancyfoot[R]{\small\textcolor{gray}{Page \thepage\ of \pageref{LastPage}}}
\renewcommand{\headrulewidth}{0.4pt}
\renewcommand{\footrulewidth}{0.4pt}

% ─── Colours ────────────────────────────────────────────────────────────────
\definecolor{mkmblue}{HTML}{1565C0}
\definecolor{codegreen}{rgb}{0,0.6,0}
\definecolor{codegray}{rgb}{0.5,0.5,0.5}
\definecolor{codepurple}{rgb}{0.58,0,0.82}
\definecolor{backcolour}{rgb}{0.95,0.95,0.92}

% ─── Hyperref setup ─────────────────────────────────────────────────────────
\hypersetup{
    colorlinks=true,
    linkcolor=mkmblue,
    citecolor=mkmblue,
    urlcolor=mkmblue
}

% ─── Code listings ──────────────────────────────────────────────────────────
\lstdefinestyle{mystyle}{
    backgroundcolor=\color{backcolour},
    commentstyle=\color{codegreen},
    keywordstyle=\color{magenta},
    numberstyle=\tiny\color{codegray},
    stringstyle=\color{codepurple},
    basicstyle=\ttfamily\footnotesize,
    breakatwhitespace=false,
    breaklines=true,
    captionpos=b,
    keepspaces=true,
    numbers=left,
    numbersep=5pt,
    showspaces=false,
    showstringspaces=false,
    showtabs=false,
    tabsize=2
}
\lstset{style=mystyle}

% ─── Section formatting ────────────────────────────────────────────────────
\setcounter{tocdepth}{3}
\setcounter{secnumdepth}{3}

% ─── Shared macros ──────────────────────────────────────────────────────────
\newcommand{\mkmtitlepage}{%
    \begin{titlepage}
    \begin{center}
    \vspace*{1cm}
    \begin{figure}[H]\centering
        \includegraphics[width=0.4\textwidth]{../shared/Logos/MKM_logo.png}
    \end{figure}
    \vspace{1.5cm}
    {\Huge\bfseries\textcolor{mkmblue}{\doctitle}\par}
    \vspace{0.5cm}
    {\Large\textcolor{gray}{\docsubtitle}\par}
    \vspace{2cm}
    {\large MKM Research Labs\par}
    \vspace{0.5cm}
    {\large Version \docversion{} --- \docdate\par}
    \vspace{0.5cm}
    {\large \docauthor\par}
    \vspace{2cm}
    \textcolor{red}{\textbf{CONFIDENTIAL --- PROPRIETARY}}
    \end{center}
    \end{titlepage}
}

\newcommand{\mkmlegalpage}{%
    \clearpage
    \thispagestyle{empty}
    \begin{center}
    \large\textbf{Legal Notice}
    \vspace{1em}

    \textbf{PROPRIETARY AND CONFIDENTIAL}
    \end{center}

    \noindent This document, including all algorithms, models, methodology, formulas,
    calculations, and intellectual concepts contained herein, constitutes the exclusive
    intellectual property and confidential trade secrets of MKM Research Labs (``MKM'').

    \vspace{0.5em}
    \noindent Any usage, reproduction, distribution, modification, or disclosure of this
    document or any portion thereof, without the express prior written authorisation from
    MKM Research Labs is strictly prohibited and constitutes an infringement of intellectual
    property rights.

    \vspace{0.5em}
    \noindent The document is provided on an ``as is'' basis. MKM Research Labs makes no
    representations or warranties of any kind, express or implied, regarding its accuracy,
    reliability, or suitability for any particular purpose.

    \vspace{1em}
    \begin{center}
    All rights reserved. \copyright\ 2021--2026 MKM Research Labs.

    \vspace{0.5em}
    Governed by the laws of the United Kingdom.
    \end{center}
    \clearpage
}

\newcommand{\mkmhistorypage}[1]{%
    \section{Document History}
    \begin{table}[ht]
        \centering
        \begin{tabular}{l|l|l|l}
            \textbf{Date} & \textbf{Version} & \textbf{Description} & \textbf{Author} \\
            \hline
            #1
        \end{tabular}
        \label{tab:revision_history}
    \end{table}
    \clearpage
}

% ─── Bank validation protocol section commands ──────────────────────────────
% Use these to ensure every model doc has the required sections.
% They are just styled \section calls — content is provided by each doc.

\newcommand{\mkmexecsummary}{\section{Executive Summary}}
\newcommand{\mkmpurpose}{\section{Model Purpose and Scope}}
\newcommand{\mkmframework}{\section{Mathematical Framework}}
\newcommand{\mkminputs}{\section{Input Parameters and Calibration}}
\newcommand{\mkmimplementation}{\section{Implementation Details}}
\newcommand{\mkmsensitivity}{\section{Sensitivity Analysis}}
\newcommand{\mkmvalidation}{\section{Validation and Backtesting}}
\newcommand{\mkmtestresults}{\subsection{Automated Test Results}\input{test_results}}
\newcommand{\mkmlimitations}{\section{Model Limitations and Known Weaknesses}}

\newcommand{\mkmgovernance}{\section{Governance Validation Questions}
\noindent The following nine questions are mandated by Chapter~5 of the
\emph{Model Risk Governance for Vendors} handbook and must be explicitly
addressed for every model in the inventory. Response status is tracked in the
Model Governance Dashboard (MRC portal). \par\medskip}

\newcommand{\mkmvqitem}[3]{\subsection{Q#1 --- #2}
\noindent\textit{``#3''}\par\medskip\noindent}

\newcommand{\mkmchangelog}{\section{Change History}}
  (before \begin{document})
%
% Required commands before \input:
%   \newcommand{\doctitle}{...}
%   \newcommand{\docsubtitle}{...}
%   \newcommand{\docversion}{1.0}
%   \newcommand{\docdate}{13-February-2026}
%   \newcommand{\docauthor}{David K Kelly}

% ─── Packages ───────────────────────────────────────────────────────────────
\usepackage[utf8]{inputenc}
\usepackage{amsmath,amsfonts,amssymb}
\usepackage{listings}
\usepackage{color,xcolor}
\usepackage{graphicx}
\usepackage{hyperref}
\usepackage[toc,page]{appendix}
\usepackage{tabto}
\usepackage{float}
\usepackage{comment}
\usepackage[en-GB]{datetime2}
\usepackage{setspace}
\usepackage{subcaption}
\usepackage{booktabs}
\usepackage{longtable}
\usepackage{adjustbox}
\usepackage{enumitem}
\usepackage{fancyhdr}
\usepackage{lastpage}

% ─── Page geometry ──────────────────────────────────────────────────────────
\onehalfspacing
\usepackage{geometry}
\geometry{a4paper, total={170mm,247mm}, left=20mm, top=25mm, bottom=25mm}

% ─── Headers and footers ───────────────────────────────────────────────────
\pagestyle{fancy}
\fancyhf{}
\fancyhead[L]{\small\textcolor{gray}{MKM Research Labs}}
\fancyhead[R]{\small\textcolor{gray}{\doctitle{} --- v\docversion}}
\fancyfoot[C]{\small\textcolor{gray}{CONFIDENTIAL}}
\fancyfoot[R]{\small\textcolor{gray}{Page \thepage\ of \pageref{LastPage}}}
\renewcommand{\headrulewidth}{0.4pt}
\renewcommand{\footrulewidth}{0.4pt}

% ─── Colours ────────────────────────────────────────────────────────────────
\definecolor{mkmblue}{HTML}{1565C0}
\definecolor{codegreen}{rgb}{0,0.6,0}
\definecolor{codegray}{rgb}{0.5,0.5,0.5}
\definecolor{codepurple}{rgb}{0.58,0,0.82}
\definecolor{backcolour}{rgb}{0.95,0.95,0.92}

% ─── Hyperref setup ─────────────────────────────────────────────────────────
\hypersetup{
    colorlinks=true,
    linkcolor=mkmblue,
    citecolor=mkmblue,
    urlcolor=mkmblue
}

% ─── Code listings ──────────────────────────────────────────────────────────
\lstdefinestyle{mystyle}{
    backgroundcolor=\color{backcolour},
    commentstyle=\color{codegreen},
    keywordstyle=\color{magenta},
    numberstyle=\tiny\color{codegray},
    stringstyle=\color{codepurple},
    basicstyle=\ttfamily\footnotesize,
    breakatwhitespace=false,
    breaklines=true,
    captionpos=b,
    keepspaces=true,
    numbers=left,
    numbersep=5pt,
    showspaces=false,
    showstringspaces=false,
    showtabs=false,
    tabsize=2
}
\lstset{style=mystyle}

% ─── Section formatting ────────────────────────────────────────────────────
\setcounter{tocdepth}{3}
\setcounter{secnumdepth}{3}

% ─── Shared macros ──────────────────────────────────────────────────────────
\newcommand{\mkmtitlepage}{%
    \begin{titlepage}
    \begin{center}
    \vspace*{1cm}
    \begin{figure}[H]\centering
        \includegraphics[width=0.4\textwidth]{../shared/Logos/MKM_logo.png}
    \end{figure}
    \vspace{1.5cm}
    {\Huge\bfseries\textcolor{mkmblue}{\doctitle}\par}
    \vspace{0.5cm}
    {\Large\textcolor{gray}{\docsubtitle}\par}
    \vspace{2cm}
    {\large MKM Research Labs\par}
    \vspace{0.5cm}
    {\large Version \docversion{} --- \docdate\par}
    \vspace{0.5cm}
    {\large \docauthor\par}
    \vspace{2cm}
    \textcolor{red}{\textbf{CONFIDENTIAL --- PROPRIETARY}}
    \end{center}
    \end{titlepage}
}

\newcommand{\mkmlegalpage}{%
    \clearpage
    \thispagestyle{empty}
    \begin{center}
    \large\textbf{Legal Notice}
    \vspace{1em}

    \textbf{PROPRIETARY AND CONFIDENTIAL}
    \end{center}

    \noindent This document, including all algorithms, models, methodology, formulas,
    calculations, and intellectual concepts contained herein, constitutes the exclusive
    intellectual property and confidential trade secrets of MKM Research Labs (``MKM'').

    \vspace{0.5em}
    \noindent Any usage, reproduction, distribution, modification, or disclosure of this
    document or any portion thereof, without the express prior written authorisation from
    MKM Research Labs is strictly prohibited and constitutes an infringement of intellectual
    property rights.

    \vspace{0.5em}
    \noindent The document is provided on an ``as is'' basis. MKM Research Labs makes no
    representations or warranties of any kind, express or implied, regarding its accuracy,
    reliability, or suitability for any particular purpose.

    \vspace{1em}
    \begin{center}
    All rights reserved. \copyright\ 2021--2026 MKM Research Labs.

    \vspace{0.5em}
    Governed by the laws of the United Kingdom.
    \end{center}
    \clearpage
}

\newcommand{\mkmhistorypage}[1]{%
    \section{Document History}
    \begin{table}[ht]
        \centering
        \begin{tabular}{l|l|l|l}
            \textbf{Date} & \textbf{Version} & \textbf{Description} & \textbf{Author} \\
            \hline
            #1
        \end{tabular}
        \label{tab:revision_history}
    \end{table}
    \clearpage
}

% ─── Bank validation protocol section commands ──────────────────────────────
% Use these to ensure every model doc has the required sections.
% They are just styled \section calls — content is provided by each doc.

\newcommand{\mkmexecsummary}{\section{Executive Summary}}
\newcommand{\mkmpurpose}{\section{Model Purpose and Scope}}
\newcommand{\mkmframework}{\section{Mathematical Framework}}
\newcommand{\mkminputs}{\section{Input Parameters and Calibration}}
\newcommand{\mkmimplementation}{\section{Implementation Details}}
\newcommand{\mkmsensitivity}{\section{Sensitivity Analysis}}
\newcommand{\mkmvalidation}{\section{Validation and Backtesting}}
\newcommand{\mkmtestresults}{\subsection{Automated Test Results}\input{test_results}}
\newcommand{\mkmlimitations}{\section{Model Limitations and Known Weaknesses}}

\newcommand{\mkmgovernance}{\section{Governance Validation Questions}
\noindent The following nine questions are mandated by Chapter~5 of the
\emph{Model Risk Governance for Vendors} handbook and must be explicitly
addressed for every model in the inventory. Response status is tracked in the
Model Governance Dashboard (MRC portal). \par\medskip}

\newcommand{\mkmvqitem}[3]{\subsection{Q#1 --- #2}
\noindent\textit{``#3''}\par\medskip\noindent}

\newcommand{\mkmchangelog}{\section{Change History}}

% Generated by docs/xva/validation_evidence.py -- do not edit.
\newcommand{\evCommit}{\texttt{04dbe701b2cf}}
\newcommand{\evGenerated}{2026-10-09T14:15:26+00:00}
\newcommand{\evPaths}{10{,}000}
\newcommand{\evConfig}{\texttt{301a1ed41a80}}
\newcommand{\evFair}{203.56}
\newcommand{\evAllPass}{all pass}


\begin{document}

\mkmtitlepage
\mkmlegalpage

\tableofcontents
\clearpage

% =========================================================================
\mkmexecsummary
% =========================================================================

MKM-XV-001 adds credit, funding and capital valuation adjustments (CVA, FVA,
KVA, with DVA disclosed) to the Physical Risk Swap. It follows Green,
\emph{XVA: Credit, Funding and Capital Valuation Adjustments} (Wiley, 2015):
every adjustment is a spread times a survival-discounted expected exposure,
and the funding strategy decides which terms exist. The model is an overlay on
the PRS pricer (MKM-PR-001): it reproduces the pricer's value at $t=0$ and does
not change it.

\paragraph{Status.} ILLUSTRATIVE, version 0.9, submitted for independent
validation. Counterparty credit, the bank's funding spread and CSA terms are
placeholders; the regulatory treatment that drives KVA is undecided. KVA is
reported beside the booked XVA and never inside it (decision D-16).

\paragraph{Findings.} On a GBP 10m, 5-year PRS with a 2\% annual trigger probability:
\begin{enumerate}[leftmargin=*]
\item Under independence, CVA and FVA are below 1bp of running spread. The
  protection buyer's CVA is almost entirely the unpaid notional during the
  settlement lag.
\item Hazard volatility and wrong-way risk move CVA by a factor of 4--13, so
  their calibration matters more than any credit input.
\item Capital dominates: KVA is about two orders of magnitude larger than CVA
  and FVA (Table~\ref{tab:ev-capital}), about 43bp of running spread against a
  fair spread of \evFair{}bp. The regulatory treatment of a PRS is worth about
  32bp of spread on this trade, and the RRAO a further 15bp.
\item Variation margin does not reach the jump to notional; initial margin
  does. A CSA is worth far more in capital than in CVA.
\end{enumerate}

\paragraph{Evidence.} All six benchmark checks pass (Table~\ref{tab:ev-checks}),
including agreement with an independent re-implementation to better than
$10^{-8}$ and the Basel Committee's own SA-CCR worked examples. Figures are
from \texttt{docs/xva/validation\_evidence.py} at commit \evCommit{}
(\evPaths{} paths, config \evConfig{}, generated \evGenerated{}): \evAllPass{}.

% =========================================================================
\mkmpurpose
% =========================================================================

\subsection{What the model answers}
What a PRS is worth to the bank once the counterparty's default, the bank's
own funding and the regulatory capital the trade consumes are priced in; and
therefore the all-in spread at which the bank should deal.

\subsection{The instrument}
A PRS is a single-name CDS with a flood trigger in place of the credit event.
The protection buyer pays a quarterly running spread in arrears while no
trigger has occurred; the seller pays the full notional (recovery 0) when the
reference gauge first passes its Severe level. The XVA model takes its
contract terms from the pricer (D-10), including first-trigger termination,
and adds a settlement lag $\Delta_s$ (provisionally 5 business days, D-09) in
which the payout is owed but unpaid.

\subsection{Scope}
\begin{itemize}[leftmargin=*]
\item In scope: gauge PRS (\texttt{TradeType = 'PRS'}), bank as buyer or
  seller, netting sets of such trades, VM and IM under a CSA, a wrong-way-risk
  stress, SA-CCR and BA-CVA capital, RRAO as a switch.
\item Out of scope: Property PRS (no term pricer, \S\ref{sec:limits}), MVA,
  TVA, warehousing (Green Ch.~14) and portfolio KVA (Ch.~15).
\end{itemize}

\subsection{Model chain}
Upstream: MKM-PR-001 (PRS value), MKM-EF-001 (event catalogue and arrival
rate). Downstream: MKM-TD-001 (Trading Desk blotter and EOD). Shares its
counterparty-capital module with MKM-CR-001 (D-07).

% =========================================================================
\mkmframework
% =========================================================================

\subsection{Green's semi-replication}
Write the all-in value as $\hat V = V + U$. Hedging the trade and both
defaults with a self-financing portfolio and applying Feynman--Kac gives
\begin{equation}
U(t) = -\int_t^T D(t,s)\,\mathbb{E}_t[f(s)]\,ds ,
\end{equation}
with one source term per adjustment:
\begin{align}
\mathrm{CVA} &= -(1-R_C)\int \lambda_C(s)\, D(0,s)\, \mathbb{E}[(V-X-\tilde I)^+]\,ds \quad(3.70, 10.31)\\
\mathrm{DVA} &= -(1-R_B)\int \lambda_B(s)\, D(0,s)\, \mathbb{E}[(V-X+I)^-]\,ds \quad(3.71)\\
\mathrm{FCA} &= -(1-R_B)\int \lambda_B(s)\, D(0,s)\, \mathbb{E}[(V-X+I)^+]\,ds \quad(9.57, 10.28)\\
\mathrm{FVA} &= -\int s_F\, D(0,s)\, \mathbb{E}[V-X]\,ds \quad(9.64\text{--}9.67)\\
\mathrm{KVA} &= -\int (\gamma_K - \phi\, r(s))\, D(0,s)\, \mathbb{E}[K(s)]\,ds \quad(13.25)
\end{align}

\subsection{Funding strategy and discounting (D-01, D-03)}
The booked strategy is Green's one-bond strategy: unilateral CVA plus a
symmetric FVA, with $D(0,s) = DF(s)\,e^{-s_F s} S_C(s)$ (discounting at
$r_F + \lambda_C$). The semi-replication view, CVA + DVA + FCA with
$D(0,s) = DF(s)\,S_B(s)\,S_C(s)$, is computed beside it and its DVA is
disclosed. The two coincide exactly when $R_B = 0$, which is a test. D-03 as
first written applied $r+\lambda_B+\lambda_C$ to every term; that is correct
only for semi-replication and was corrected in WP1.

\subsection{Why a PRS is different}
Exposure is a jump, not a diffusion: before a trigger the value is a small
spread MTM; after it, $\pm N$ until settlement. Under deterministic hazard the
pre-trigger value is deterministic, so
\begin{equation}
\mathbb{E}[V(s)^+] = S_E(s)\,V(s)^+ + \int_{\max(0,s-\Delta_s)}^{\min(s,T)} \lambda_E S_E(u)\, P(u,s)^+\,du,
\end{equation}
with $P(u,s)$ the pending payout. This is exact, and is the closed form the
Monte Carlo must reproduce.

\subsection{Stochastic flood hazard: CIR++ (D-15)}
A catchment factor $Y$ follows CIR with $Y(0)=\theta=1$:
$dY = \kappa(1-Y)\,dt + \sigma\sqrt{Y}\,dW$. Gauge $g$'s trigger intensity is
$\lambda_g Y(t) + \varphi_g(t)$, with the shift fitted so the model's survival
at $t=0$ is exactly the pricer's (Brigo--Mercurio). Given $Y(t)=y$ and no
trigger,
\begin{equation}
S(t,u) = e^{-(\Phi_g(u)-\Phi_g(t))}\,A(u-t)\,e^{-B(u-t)\lambda_g y},
\end{equation}
so $V(t)$ is closed form on every path. $A$ is evaluated in a form that is
stable as $\sigma\to0$ (the textbook form loses precision there, which the
tests caught). $Y$ is stepped exactly by its noncentral $\chi^2$ transition.

\subsection{Correlated triggers}
The MKM-EF-001 catalogue supplies dependence: catchment events arrive as a Cox
process in the factor clock $C(t)=\int_0^t Y$, are drawn by catalogue weight,
and their Severe flags say which gauges they trigger. Each gauge's marginal is
held to the market rate by thinning (catalogue rate above market) or a
$Y$-driven top-up (below).

\subsection{Collateral, netting and wrong-way risk}
Variation margin follows Green's lookback (6.39) with threshold and minimum
transfer amount over the margin period of risk; initial margin is static in
each direction. Without close-out netting, exposures add trade by trade. The
wrong-way-risk stress weights each path by $e^{\beta Z}/\mathbb{E}[e^{\beta Z}]$,
$Z$ the number of the netting set's gauges triggered so far, normalised so
each period's default probability is unchanged (Green 7.17).

\subsection{Capital (D-06, D-07, D-16)}
SA-CCR EAD and BA-CVA (reduced) are computed in every Monte Carlo state
(Green \S13.4). How a PRS enters SA-CCR is a switch: commodity ``other''
(SF 18\%) or a single-name credit analogy. Capital held is RWA times the
bank's 10.5\% ratio, shared with MKM-CR-001.

% =========================================================================
\mkminputs
% =========================================================================

\begin{longtable}{p{4.2cm}p{2.6cm}p{7.2cm}}
\toprule
\textbf{Parameter} & \textbf{Value} & \textbf{Source / status} \\
\midrule
Flood trigger hazard & market state & Pricer's gauge rate (MKM-PR-001 / EF-001) \\
Risk-free curve & market state & Desk yield curve (D-13) \\
Counterparty hazard & rating proxy & \texttt{RATING\_CDS\_SPREADS}: \textbf{placeholder} (D-11) \\
$R_C$, $R_B$ & 40\%, 40\% & Market convention \\
$s_F$ & 75bp & \textbf{Placeholder}, Treasury input (D-12) \\
$\kappa$, $\sigma$ & 0.5, 0.3 & \textbf{Placeholder}, no calibration owner (D-15) \\
$\Delta_s$ & 5 business days & \textbf{Provisional}, Product/Legal (D-09) \\
$\beta$ (WWR) & 0 & Stress only (D-05) \\
$\gamma_K$, $\phi$ & 12\%, 0 & Shared with MKM-CR-001 (D-06) \\
Capital ratio & 10.5\% & Shared with MKM-CR-001 \\
SA-CCR treatment & commodity other & Regulatory paper Q2 (D-16) \\
RRAO & off & Regulatory paper Q3 \\
Paths, seed & 10{,}000, 20261009 & D-08 \\
Time cells & 26 per quarter & D-14 \\
\bottomrule
\end{longtable}

% =========================================================================
\mkmimplementation
% =========================================================================

{\small
\begin{longtable}{p{6.4cm}p{8.4cm}}
\toprule
\textbf{Module} & \textbf{Role} \\
\midrule
\texttt{config/xva.py} & Every parameter; frozen, fingerprinted \\
\texttt{models/xva/curves.py} & Pricer-identical discounting; hazard curves; CDS bootstrap \\
\texttt{models/xva/valuation.py} & $V(t)$, fair spread, pending payout \\
\texttt{models/xva/exposure\_analytic.py} & Exact EPE/ENE/EE \\
\texttt{models/xva/integrals.py} & Kernels and integrals; the only place discounting lives \\
\texttt{models/xva/analytic.py} & WP1 engine; exact break-even spread \\
\texttt{models/xva/\{cir,events,paths\}.py} & CIR++, correlated triggers, grid and path values \\
\texttt{models/xva/netting.py} & CSA and netting sets from the CDM \\
\texttt{models/xva/monte\_carlo.py} & WP2 engine; per-path estimates with SE; Euler allocation \\
\texttt{models/capital/counterparty.py} & SA-CCR, BA-CVA, risk weights (shared, D-07) \\
\texttt{models/xva/capital.py} & Capital per Monte Carlo state \\
\texttt{models/xva/desk.py},\newline\texttt{routes/trading/xva.py} & Trading Desk integration (WP4) \\
\bottomrule
\end{longtable}}

A 16-trade netting set runs in about 10 seconds at 10{,}000 paths; memory is
$O(\text{paths}\times\text{nodes})$ whatever the number of trades. The desk
runs XVA on demand, at EOD and for quotes; the blotter reads the stored run
(D-17).

% =========================================================================
\mkmvalidation
% =========================================================================

\subsection{Benchmarks}
% Generated by docs/xva/validation_evidence.py -- do not edit.
\begin{table}[H]\centering\small
\caption{Benchmark checks}\label{tab:ev-checks}
\begin{tabular}{lrrl}
\toprule
Check & Result & Tolerance & Outcome \\
\midrule
Pricer reconciliation (fair spread, annuity) & $4.2\times10^{-16}$ & $<10^{-12}$ & pass \\
Independent re-implementation, all terms & $3.7\times10^{-11}$ & $<10^{-8}$ & pass \\
Monte Carlo against analytic, max $|z|$ & 1.56 & $<3$ & pass \\
Time rule error at 26 cells / quarter & $4.6\times10^{-4}$ & $<5\times10^{-4}$ & pass \\
CIR++ fit to the market curve & $2.2\times10^{-16}$ & $<10^{-12}$ & pass \\
BCBS 279 Example 2 / 3 EAD & 381 / 5406 & 381 / 5{,}406 & pass \\
\bottomrule
\end{tabular}
\end{table}


The independent re-implementation (\texttt{tests/models/xva/\_reference.py})
shares no code with the engine and integrates adaptively. Green's two
derivations of CVA, by expectation and by replication, agree within four
standard errors when default times are sampled directly.

\subsection{Monte Carlo against the closed form}
% Generated by docs/xva/validation_evidence.py -- do not edit.
\begin{table}[H]\centering\small
\caption{Monte Carlo ($\sigma=0$) against the analytic engine}\label{tab:ev-mc-recon}
\begin{tabular}{lllrrrr}
\toprule
Bank & View & Term & MC & SE & Analytic & $z$ \\
\midrule
buys & one-bond & CVA & $-$1{,}974 & 4.7 & $-$1{,}967 & $-1.31$ \\
buys & one-bond & FVA & $-$1{,}154 & 3.2 & $-$1{,}150 & $-1.34$ \\
buys & one-bond & TOTAL & $-$3{,}128 & 7.9 & $-$3{,}117 & $-1.32$ \\
buys & semi-replication & CVA & $-$1{,}960 & 4.7 & $-$1{,}953 & $-1.31$ \\
buys & semi-replication & DVA & 78 & 0.2 & 78 & $-1.56$ \\
buys & semi-replication & FCA & $-$1{,}225 & 2.9 & $-$1{,}221 & $-1.31$ \\
buys & semi-replication & TOTAL & $-$3{,}107 & 7.8 & $-$3{,}096 & $-1.32$ \\
sells & one-bond & CVA & $-$127 & 0.4 & $-$128 & $+1.56$ \\
sells & one-bond & FVA & 1{,}154 & 3.2 & 1{,}150 & $+1.34$ \\
sells & one-bond & TOTAL & 1{,}027 & 3.6 & 1{,}022 & $+1.37$ \\
sells & semi-replication & CVA & $-$124 & 0.4 & $-$125 & $+1.56$ \\
sells & semi-replication & DVA & 1{,}225 & 2.9 & 1{,}221 & $+1.31$ \\
sells & semi-replication & FCA & $-$78 & 0.2 & $-$78 & $+1.56$ \\
sells & semi-replication & TOTAL & 1{,}023 & 3.5 & 1{,}018 & $+1.37$ \\
\bottomrule
\end{tabular}
\end{table}


\subsection{Results}
% Generated by docs/xva/validation_evidence.py -- do not edit.
\begin{table}[H]\centering\small
\caption{Analytic engine: GBP 10m, 5y, 2\% trigger, at the fair spread}\label{tab:ev-analytic}
\begin{tabular}{llrrrrrr}
\toprule
Ctpy & Bank & CVA & FVA & Booked & DVA & Semi-rep & bp \\
\midrule
A & buys & $-$91 & 663 & 573 & 778 & 566 & +0.14 \\
A & sells & $-$574 & $-$663 & $-$1{,}237 & 123 & $-$1{,}223 & +0.30 \\
BBB & buys & $-$192 & 647 & 454 & 759 & 449 & +0.11 \\
BBB & sells & $-$1{,}219 & $-$647 & $-$1{,}866 & 120 & $-$1{,}844 & +0.46 \\
BB & buys & $-$448 & 603 & 156 & 708 & 154 & +0.04 \\
BB & sells & $-$2{,}834 & $-$603 & $-$3{,}438 & 112 & $-$3{,}399 & +0.88 \\
\bottomrule
\end{tabular}
\end{table}
\begin{table}[H]\centering\small
\caption{Monte Carlo engine, BBB-proxy counterparty}\label{tab:ev-mc}
\begin{tabular}{lrrr}
\toprule
Case & CVA $\pm$ SE & FVA & Booked \\
\midrule
buys, sigma 0 & $-$201 $\pm$ 6 & 639 & 437 \\
buys, sigma 0.3 & $-$933 $\pm$ 11 & 653 & $-$279 \\
buys, WWR beta 3 & $-$2{,}677 $\pm$ 75 & 653 & $-$2{,}024 \\
buys, VM & $-$622 $\pm$ 5 & 13 & $-$610 \\
buys, VM + IM = notional & $-$1 $\pm$ 0 & 13 & 12 \\
sells, sigma 0.3 & $-$1{,}970 $\pm$ 13 & $-$653 & $-$2{,}623 \\
sells, WWR beta 3 & $-$1{,}183 $\pm$ 8 & $-$653 & $-$1{,}837 \\
sells, VM & $-$643 $\pm$ 5 & $-$13 & $-$656 \\
\bottomrule
\end{tabular}
\end{table}
\begin{table}[H]\centering\small
\caption{Capital and KVA, bank buys protection; bp is $|$KVA$|/(NA)$ of running spread}\label{tab:ev-capital}
\begin{tabular}{lrrrr}
\toprule
Treatment & EAD today & Capital today & KVA $\pm$ SE & bp \\
\midrule
commodity other & 2{,}520{,}000 & 538{,}129 & $-$183{,}869 $\pm$ 267 & 43.08 \\
credit analogy (ref. B) & 990{,}972 & 211{,}616 & $-$46{,}463 $\pm$ 61 & 10.89 \\
commodity + RRAO & 2{,}520{,}000 & 669{,}379 & $-$247{,}527 $\pm$ 373 & 58.00 \\
commodity, VM & 756{,}000 & 161{,}439 & $-$57{,}677 $\pm$ 88 & 13.52 \\
\bottomrule
\end{tabular}
\end{table}


\subsection{Automated tests}
{\small
\begin{longtable}{p{7.6cm}p{7.2cm}}
\toprule
\textbf{Suite} & \textbf{Purpose} \\
\midrule
\texttt{tests/models/xva} & Curves, valuation, benchmark, limiting cases, CIR, events, netting, Monte Carlo, capital, desk \\
\texttt{tests/models/capital/test\_counterparty.py} & BCBS 279 Annex 4a examples \\
\texttt{tests/routes/trading/test\_xva\_routes.py} & Run, blotter, quote, EOD \\
\texttt{tests/js/blotter-table.test.js} & XVA and KVA columns \\
\texttt{tests/port/cdm/test\_credit\_support.py} & Netting and CSA CDM fields \\
\bottomrule
\end{longtable}}

\noindent Coverage of \texttt{models.xva} and \texttt{models.capital} is 100\%.

\subsection{Outcomes analysis}
Triggers are rare, so exposures cannot be backtested meaningfully. Outcomes
analysis rests on the MKM-EF-001 hazard backtests, the daily valuation KPI
(engine $V(0)$ against the front-office MTM, 5bp of notional) and the
$\Delta$XVA P\&L line at EOD.

% =========================================================================
\mkmsensitivity
% =========================================================================

% Generated by docs/xva/validation_evidence.py -- do not edit.
\begin{table}[H]\centering\small
\caption{Sensitivities: booked XVA and KVA (GBP)}\label{tab:ev-sens}
\begin{tabular}{lrrrr}
\toprule
Case & Buys XVA & Buys KVA & Sells XVA & Sells KVA \\
\midrule
base & $-$279 & $-$183{,}869 & $-$2{,}623 & $-$187{,}668 \\
counterparty hazard x0.5 & 193 & $-$187{,}359 & $-$1{,}676 & $-$191{,}236 \\
counterparty hazard x2 & $-$1{,}166 & $-$177{,}169 & $-$4{,}394 & $-$180{,}820 \\
R\_C 20\% & $-$590 & $-$183{,}869 & $-$3{,}280 & $-$187{,}668 \\
R\_C 60\% & 32 & $-$183{,}869 & $-$1{,}967 & $-$187{,}668 \\
s\_F 0bp & $-$947 & $-$186{,}498 & $-$2{,}003 & $-$190{,}355 \\
s\_F 150bp & 366 & $-$181{,}294 & $-$3{,}221 & $-$185{,}036 \\
sigma 0 & 437 & $-$182{,}235 & $-$1{,}855 & $-$186{,}002 \\
sigma 0.6 & $-$1{,}324 & $-$184{,}497 & $-$3{,}762 & $-$188{,}455 \\
kappa 0.25 & $-$887 & $-$184{,}091 & $-$3{,}267 & $-$187{,}941 \\
kappa 1.0 & 221 & $-$183{,}108 & $-$2{,}121 & $-$186{,}910 \\
WWR beta 1 & $-$523 & $-$183{,}869 & $-$2{,}495 & $-$187{,}668 \\
settlement 0d & 31 & $-$183{,}185 & $-$2{,}738 & $-$187{,}193 \\
settlement 20d & $-$1{,}239 & $-$183{,}958 & $-$2{,}248 & $-$187{,}103 \\
credit-analogy SA-CCR & $-$279 & $-$46{,}463 & $-$2{,}623 & $-$50{,}162 \\
RRAO on & $-$279 & $-$247{,}527 & $-$2{,}623 & $-$251{,}326 \\
gamma\_K 8\% & $-$279 & $-$122{,}579 & $-$2{,}623 & $-$125{,}112 \\
gamma\_K 15\% & $-$279 & $-$229{,}836 & $-$2{,}623 & $-$234{,}585 \\
phi = 1 & $-$279 & $-$120{,}294 & $-$2{,}623 & $-$122{,}805 \\
\bottomrule
\end{tabular}
\end{table}


The base case is the trade above at the fair spread against a BBB-proxy
counterparty, $\sigma=0.3$, commodity SA-CCR. KVA is the larger figure in
every row; among XVA inputs, hazard volatility and the settlement lag move the
buyer most, counterparty hazard the seller.

% =========================================================================
\mkmlimitations\label{sec:limits}
% =========================================================================

\begin{longtable}{p{1cm}p{13.5cm}}
\toprule
L-1 & Classical MPoR freezes collateral but not trade cash flows: a buyer that posted VM is briefly over-collateralised when it pays premium (Andersen--Pykhtin--Sokol). \\
L-2 & The event clock integrates $Y$ by the trapezoid rule; exact at $\sigma=0$. \\
L-3 & WWR driver is the netting set's triggered gauges, not the counterparty's real flood book. \\
L-4 & Euler allocation is exact only for uncollateralised netted sets; residuals reported. \\
L-5 & CIR++ is fitted to a flat trigger hazard (the pricer's convention). \\
L-6 & One seed, no variance reduction; the closed form is available as a control variate. \\
L-7 & Every regulatory treatment is provisional; KVA is not booked. \\
L-8 & The credit analogy needs a ``reference rating'' for a flood trigger. \\
L-9 & BA-CVA applies CRE32's one-year maturity floor; to be confirmed. \\
L-10 & Standalone BA-CVA: no diversification across counterparties. \\
L-11 & No market-risk capital beyond the RRAO switch. \\
L-12 & $\gamma_K$ is a single constant. \\
L-13 & Desk uses one flat hazard per gauge; the front office's term structure is not used. \\
L-14 & KVA allocated to trades pro rata to notional. \\
L-15 & Netting set = counterparty (first CDM netting set). \\
L-16 & No quote panel or ``Run XVA'' button in the UI yet. \\
L-17 & Integration not yet exercised on the live book. \\
\bottomrule
\end{longtable}

\paragraph{Findings against other models.} F-1: the pricer's MTM uses a fair
spread floored at 2bp. F-2: the pricer ignores the settlement lag. F-3: the
pricer rounds MTM to the penny. The trade record carries a 6M premium
frequency while the pricer prices quarterly (D-10).

% =========================================================================
\mkmgovernance
% =========================================================================

\mkmvqitem{1}{Ronseal Test}{Does the model do what it says on the tin?}
Yes for the scope above. It reproduces the pricer at $t=0$ and Green's
integrals to $10^{-8}$; its Monte Carlo reproduces the closed form.

\mkmvqitem{2}{Accuracy and Stability}{Does the model provide timely, accurate and stable results?}
Accurate against every benchmark. Every Monte Carlo figure carries its
standard error; results are reproducible from the seed. Runs take seconds per
netting set.

\mkmvqitem{3}{Resource Usage}{Does the model consume excessive computational resources?}
No: about 1s per trade at 10{,}000 paths, memory independent of trade count.
It is kept out of per-request revaluation for this reason.

\mkmvqitem{4}{Override Controls}{Are there controls over manual overrides?}
No manual overrides exist. Every parameter is in \texttt{config/xva.py},
fingerprinted into each run's audit block.

\mkmvqitem{5}{Output Consumption}{How is the output consumed for risk decision-making?}
Blotter XVA and KVA memo columns, EOD $\Delta$XVA, and incremental quotes.
Nothing reaches a price or P\&L until validation, and KVA not until the
regulatory paper is answered.

\mkmvqitem{6}{Model Chaining}{Is the model's output used as input to another model?}
Yes: it feeds MKM-TD-001, and is fed by MKM-PR-001 and MKM-EF-001; it shares
capital code with MKM-CR-001.

\mkmvqitem{7}{Weaknesses}{Does the model have weaknesses and limitations of use?}
Yes: \S\ref{sec:limits}. The material ones are placeholder inputs (credit,
funding, $\sigma$), the open regulatory treatment, and L-1.

\mkmvqitem{8}{Monitoring}{What type of monitoring is required?}
Daily valuation KPI; standard error per component; $\Delta$XVA P\&L explain;
carry artefacts (grid roll, seed); annual recalibration of $\sigma$, $\kappa$.

\mkmvqitem{9}{Overall Risk}{Is the overall model risk acceptable?}
Not yet for use: Tier 1, Amber. Acceptable for validation, with inputs and
the regulatory treatment as gating items.

% =========================================================================
\mkmhistorypage{9-Oct-2026 & 0.9 & Built (WP0--WP4); submitted for validation & David K Kelly \\}
% =========================================================================

\end{document}
